% Generated TeX body; no PDF or document-style claim.
% Maintained source: manuscripts/matching-and-routing/rank-sensitive-layer-count.md
% Recovered-source SHA-256: ad1ec885f88ccf5b11a826186910cfed5ab067764d7a69340f192e3a795a7aaf
% Rebuild: python3 tools/build_manuscripts.py
% Common packages: amsmath, amssymb, hyperref, longtable, booktabs, array.
% Shared original Subset Sum preamble: manuscripts/paper-style.tex.

\hypertarget{af-6a8212874a-rank-sensitive-counting-for-certified-haxell-layers}{%
\section{Rank-sensitive counting for certified Haxell layers}\label{af-6a8212874a-rank-sensitive-counting-for-certified-haxell-layers}}

\textbf{Further analysis, 25 October 2026.} This note strengthens the weighted count from \href{https://github.com/icekylinx/algorithm-frontier/blob/main/problems/strong-haxell-matching/proof.md}{strong haxell matching proof}. It applies both to the no-superposition algorithm and, through the constant number of adopted-reset epochs, to the original BHHS operations reanalyzed in \href{https://github.com/icekylinx/algorithm-frontier/blob/main/manuscripts/matching-and-routing/public-hypergraph-algorithm-weighted-reanalysis.md}{public-hypergraph-algorithm-weighted-reanalysis.md}. It does not alter the previously delivered checkpoint archive.

The improvement uses the cutoff at which a parent is still allowed to create another child. The resulting number of children depends on the hyperedge rank, rather than on the rank times the degree cap. The root contributes a separate logarithmic factor, which cannot be removed by a simple active-incidence accounting argument.

\hypertarget{af-6a8212874a-statement}{%
\subsection{1. Statement}\label{af-6a8212874a-statement}}

Use the common explicit-hypergraph convention: \texttt{q=\textbar{}A\textbar{}} left vertices, at most \texttt{r} right vertices per edge, input volume \texttt{p=sum\_e\ \textbar{}e\textbar{}}, and degree cap \texttt{Delta\textgreater{}=4}. Write

\[
U=\log(p+4),\quad h=\log(U+4),\quad s=\log\Delta,
 \qquad
 a=\min\{\log(2+s),\log(2+\log(r+2))\}.
 \tag{1}
\]

After the standard trivial checks, \texttt{Delta\ r\^{}2\textless{}=p}, so \texttt{s\textless{}=U}. Assume either the strong condition

\[
\tau(E_S)\ge\Delta r^2 |S|
                    \quad(S\subseteq A),
 \tag{2}
\]

or use the explicit growth-test/certificate alternative of the no-superposition algorithm. In the latter case the count applies until a violation is returned.

The number of layer builds is at most

\[
O\!\left(\min\{\log(q+1),\,1+s/a\}\right)
        \exp O\!\left(\frac{a\log(q+1)}s\right).
 \tag{3}
\]

Consequently the explicit implementation takes

\[
O(|A|+|B|)+p\exp O\!\left(
             h+\frac{a\log(q+1)}s\right)
 \tag{4}
\]

time. The original source's operations have the same bound under (2); successful superpositions add only a constant number of epochs. The faster no-promise certifying outcome is stated for the explicitly tested implementation.

For constant rank, (3) specializes to

\[
O\!\left(\min\{\log(q+1),1+\log\Delta\}\right)
      \exp O\!\left(\frac{\log(q+1)}{\log\Delta}\right).
 \tag{5}
\]

In particular, for \texttt{Delta\textgreater{}=p\^{}epsilon}, this gives

\[
O(|A|+|B|)+p\log(p+4)\exp O(1/\epsilon)
 \tag{6}
\]

for constant \texttt{r}. The constants in this expression are uniform in \texttt{epsilon}; its feasible parameter range is inherited from \texttt{Delta\ r\^{}2\textless{}=p}.

\hypertarget{af-6a8212874a-the-child-cutoff-removes-the-factor-delta-from-the-fanout}{%
\subsection{2. The child cutoff removes the factor Delta from the fanout}\label{af-6a8212874a-the-child-cutoff-removes-the-factor-delta-from-the-fanout}}

Fix a positive-index layer instance \texttt{i}, and let \texttt{W=\textbar{}Y\_\textless{}i\textbar{}\textgreater{}0}. This lower prefix remains fixed during the layer's lifetime. The selected edges obey their cap, so

\[
|X_i|\le\Delta|Y_{i-1}|\le\Delta W.
\]

Every selected edge has at most \texttt{r} matching blockers, hence

\[
|Y_i|\le r\Delta W.
 \tag{7}
\]

In the no-superposition algorithm, \texttt{Y\_i} only decreases. In the public algorithm, (7) holds at the start of each epoch and throughout the epoch; an adopted superposition starts a new one.

A direct child is born only when the parent is noncollapsible. The growth and noncollapse inequalities give

\[
z=W+|Y_i|\ge\kappa W,
 \qquad\kappa=1+9\Delta/100.
 \tag{8}
\]

Within one epoch, the prefix sizes of successive children satisfy \texttt{z\_(j+1)\textless{}=rho\ z\_j}, where \texttt{rho=.99}. This follows from the child's swap: it removes at least one percent of the then-current \texttt{W+\textbar{}Y\_i\textbar{}}. Therefore all child birth sizes in the epoch lie in the interval

\[
[\kappa W,(1+r\Delta)W].
\]

The ratio of its endpoints is bounded independently of \texttt{Delta}:

\[
\frac{1+r\Delta}{1+9\Delta/100}
                   \le\frac{100r}{9}<12r.
 \tag{9}
\]

An epoch consequently has at most

\[
L_r=2+\left\lceil\frac{\log(12r)}{-\log\rho}\right\rceil
        =O(\log(r+2))
 \tag{10}
\]

created children. The possible last living child is included: the argument applies to the list of births already present in any finite execution prefix, without assuming it will terminate.

For the no-superposition algorithm there is one epoch. For the public algorithm there are at most \texttt{J=26}, because the selected count starts above \texttt{Delta\ W/10}, is bounded by \texttt{Delta\ W}, and each adopted superposition increases it by a factor at least \texttt{1.1}. Thus a positive-index instance has at most \texttt{J\ L\_r} direct children.

\hypertarget{af-6a8212874a-rank-sensitive-contracting-weights}{%
\subsection{3. Rank-sensitive contracting weights}\label{af-6a8212874a-rank-sensitive-contracting-weights}}

Choose

\[
\lambda_r=\frac{\log(2J L_r)}{\log\kappa}.
 \tag{11}
\]

Give the parent instance weight \texttt{W\^{}(-lambda\_r)}. Every child has prefix size at least \texttt{kappa\ W} by (8), so

\[
\sum_{child} W_{child}^{-\lambda_r}
       \le J L_r\kappa^{-\lambda_r} W^{-\lambda_r}
       =\tfrac12W^{-\lambda_r}.
 \tag{12}
\]

For \texttt{Delta\textgreater{}=4}, one has \texttt{log\ kappa=Theta(log\ Delta)}; for example \texttt{kappa\textgreater{}=Delta\^{}(1/5)}. Equation (10) therefore implies

\[
\lambda_r=\Theta\!\left(
                   \frac{\log(2+\log(r+2))}s\right).
\]

Independently, the rank-free geometric-series proof supplies \texttt{lambda\_u=C\ log(2+s)/s} with the same one-half child-weight contraction, including \texttt{J=26} epochs, for an absolute constant \texttt{C}. Take \texttt{lambda=min\{lambda\_r,lambda\_u\}} and use the corresponding proof. Thus

\[
\lambda=\Theta(a/s).
 \tag{13}
\]

The weights are analysis objects only. No real arithmetic or evaluation of powers is required by the algorithm.

\hypertarget{af-6a8212874a-the-root-factor-must-be-retained}{%
\subsection{4. The root factor must be retained}\label{af-6a8212874a-the-root-factor-must-be-retained}}

The root's children have prefix sizes equal to successive positive unmatched counts. These counts decrease by a factor at most \texttt{rho} between children. They are at least one, so their weights sum to at most

\[
\min\left\{O(\log(q+1)),\frac1{1-\rho^\lambda}\right\}
       =O\!\left(\min\{\log(q+1),1+1/\lambda\}\right).
 \tag{14}
\]

The first bound uses the number of root children and the fact that every weight is at most one. The second reverses their geometrically decreasing prefix-size sequence.

Summing (12) over the forest of positive-index layer instances shows that total node weight is at most twice (14). Every node weight is at least \texttt{q\^{}(-lambda)}. Equations (13)--(14) prove (3), in every finite execution prefix. They also prove termination: another unmatched left vertex would require either another finite build or a returned violation certificate.

When \texttt{r} is constant and \texttt{s} is comparable to \texttt{U}, the main exponent \texttt{lambda\ log\ q} can be constant, but the prefactor in (14) can grow like \texttt{log\ p}. It cannot generally be absorbed into that constant exponent. Keeping \texttt{h} in (4) is essential for this proof.

\hypertarget{af-6a8212874a-concrete-constant-rank-regimes}{%
\subsection{5. Concrete constant-rank regimes}\label{af-6a8212874a-concrete-constant-rank-regimes}}

The following table gives overheads multiplying \texttt{p}, with the same strong promise (2). It omits the one-time unused vertex-list reading cost.

\begin{longtable}[]{@{}
  >{\raggedright\arraybackslash}p{(\columnwidth - 4\tabcolsep) * \real{0.2727}}
  >{\raggedleft\arraybackslash}p{(\columnwidth - 4\tabcolsep) * \real{0.3636}}
  >{\raggedleft\arraybackslash}p{(\columnwidth - 4\tabcolsep) * \real{0.3636}}@{}}
\toprule\noalign{}
\begin{minipage}[b]{\linewidth}\raggedright
Strength parameter
\end{minipage} & \begin{minipage}[b]{\linewidth}\raggedleft
Previous rank-free weighted bound
\end{minipage} & \begin{minipage}[b]{\linewidth}\raggedleft
Rank-sensitive bound for constant \texttt{r}
\end{minipage} \\
\midrule\noalign{}
\endhead
\bottomrule\noalign{}
\endlastfoot
\texttt{Delta\textgreater{}=p\^{}epsilon} & \texttt{(log\ p)\^{}\{O(1/epsilon)\}} & \texttt{log(p+4)\ exp\ O(1/epsilon)} \\
\texttt{Delta\textgreater{}=exp(U\^{}beta)}, fixed \texttt{0\textless{}beta\textless{}1} & \texttt{exp\ O(U\^{}(1-beta)\ log\ U)} & \texttt{U\ exp\ O(U\^{}(1-beta))} \\
\texttt{Delta\textgreater{}=U\^{}c}, fixed \texttt{c\textgreater{}0} & \texttt{exp\ O(U\ log\ log\ U/log\ U)} & \texttt{U\ exp\ O(U/log\ U)} \\
\end{longtable}

These are bounds for the layer method, not lower bounds on the cost of the task. Take their minimum with the elementary linear-time regimes below. In particular the first row is not an improvement over elementary greedy matching when \texttt{Delta\textgreater{}=sqrt(p)}; for fixed \texttt{epsilon\textgreater{}=1/2} that entire parameter range is already linear after input reading.

For nonconstant rank, formula (1) selects the better of its rank-sensitive and rank-free dependence. Both rates also follow for the public BHHS operation sequence by the constant-epoch argument. This is a further quantitative analysis, not an exclusive advantage of a different explicit-hypergraph algorithm.

\hypertarget{af-6a8212874a-an-elementary-linear-time-regime}{%
\subsubsection{An elementary linear-time regime}\label{af-6a8212874a-an-elementary-linear-time-regime}}

There is a useful separate endpoint that does not require layers. If

\[
\Delta r>q-1,
 \tag{14a}
\]

scan the incident edge lists of the left vertices one at a time, and choose any edge avoiding the right resources already used by the matching. Before the last choice, those resources number at most \(r(q-1)<\Delta r^2\). The singleton case of (2) guarantees a permitted edge. Every edge list is scanned at most once, so this takes \(O(|A|+|B|+p)\) time. Without the promise, a failed left vertex and the matching's current right resources are an immediate singleton strong-Haxell violation certificate.

One should first test whether \(|B_{\rm used}|<\Delta r^2q\), in which case all left vertices and all used right resources are a valid violation. If that test passes and \(\Delta^2r^3\ge |B_{\rm used}|\), then

\[
q\le\frac{|B_{\rm used}|}{\Delta r^2}
   \le\Delta r,
\]

so (14a) holds. The direct input condition \(q\le\Delta r\) already suffices, without the resource-count inference. Since \(|B_{\rm used}|\le p\), the more convenient condition \(\Delta^2r^3\ge p\), and in particular \(\Delta\ge\sqrt p\), also suffices for this elementary linear-time certifying endpoint.

Another elementary preprocessing observation applies throughout the parameter range. Scan all left vertices greedily once, producing a maximal matching \(M\) and the uncovered left set \(U_0\). Every edge incident to \(U_0\) meets \(B(M)\), because the matching only grew after a failed scan. Under (2),

\[
\Delta r^2|U_0|
   \le\tau(E_{U_0})\le |B(M)|\le r|M|,
 \qquad
 |U_0|\le\frac{q}{\Delta r+1}.
 \tag{14b}
\]

Without the promise, if \(|B(M)|<\Delta r^2|U_0|\), then \((U_0,B(M))\) is itself a valid strong-Haxell violation. This gives a useful small uncovered set, but does not automatically solve that set recursively: its available edges may still meet the existing matching. An augmenting or reassignment argument is still needed.

These endpoints do not improve the root-prefactor proof itself. They are separate elementary implementations or preprocessing observations. The example below is explicitly an obstruction for the ordinary layer-list scan. A preliminary greedy matching can avoid it, so the example is not an obstruction for this or every other algorithm.

\hypertarget{af-6a8212874a-rank-one-already-has-a-stronger-classical-phase-bound}{%
\subsubsection{Rank one already has a stronger classical phase bound}\label{af-6a8212874a-rank-one-already-has-a-stronger-classical-phase-bound}}

When \texttt{r=1}, the hypergraph is a bipartite graph, and (2) is simply \(|N(S)|\ge\Delta|S|\). In this case use the classical \href{https://doi.org/10.1137/0202019}{Hopcroft--Karp algorithm} and its linear-time phases of maximal collections of shortest augmenting paths.

For any nonperfect matching, start alternating BFS from an unmatched left vertex. If the currently reachable left set \(S_j\) has no free right neighbor, every vertex of \(N(S_j)\) is matched and its distinct left mate is reachable in the next step. A matched edge incident to an already reached left vertex only gives a mate that was already reached, so it does not invalidate this counting. Therefore

\[
|S_{j+1}|\ge|N(S_j)|\ge\Delta|S_j|.
\]

Since there are only \(q\) left vertices, an augmenting path has length \(O(1+\log(q+1)/\log\Delta)\). The shortest augmenting length strictly increases after a Hopcroft--Karp phase, so the total time is

\[
O(|A|+|B|)+O\!\left(p\left(1+
                       \frac{\log(q+1)}{\log\Delta}\right)\right).
 \tag{14c}
\]

In particular this is \(O(p(1+1/\epsilon))\), plus input reading, when \(\Delta\ge p^\epsilon\). The new layer analysis is not a speed improvement in this graph case. Its nontrivial constant-rank comparison concerns \(r\ge2\) outside the elementary greedy regime.

\hypertarget{af-6a8212874a-why-scanning-only-active-incidences-does-not-eliminate-the-root-logarithm}{%
\subsection{6. Why scanning only active incidences does not eliminate the root logarithm}\label{af-6a8212874a-why-scanning-only-active-incidences-does-not-eliminate-the-root-logarithm}}

Here is an explicit family satisfying the strong promise, of constant rank two, on which the direct edge-list implementation spends \texttt{Omega(p\ log\ p)} time. The same lower estimate applies if each build scans only edges incident to its active left vertices. This is an implementation example, not a lower bound for every possible algorithm or data structure.

\hypertarget{af-6a8212874a-a-3-uniform-example-with-large-delta}{%
\subsubsection{A 3-uniform example with large Delta}\label{af-6a8212874a-a-3-uniform-example-with-large-delta}}

Take \texttt{q} left vertices \texttt{a\_1,...,a\_q}, let \texttt{Delta=q\^{}3}, and create original right resources \texttt{x\_1,...,x\_(4Delta\ q)}. For every \texttt{j\textgreater{}=2} and \texttt{1\textless{}=t\textless{}=4Delta\ j}, introduce a distinct private right resource \texttt{b\_(j,t)} and the hyperedge

\[
\{a_j,x_t,b_{j,t}\}.
\]

For \texttt{a\_1}, include all hyperedges

\[
\{a_1,x_u,x_v\},
             \qquad1\le u<v\le4\Delta+1.
\]

Every edge has exactly two right vertices and total size three. There are no duplicate hyperedges. The input volume is

\[
p=3\binom{4\Delta+1}{2}
       +3\sum_{j=2}^q4\Delta j
       =\Theta(\Delta^2+\Delta q^2)=\Theta(q^6).
 \tag{15}
\]

It is dominated by the edges incident to \texttt{a\_1}.

For \texttt{S=\{a\_1\}}, the right-side hitting number is the vertex cover number of a clique on \texttt{4Delta+1} resources, namely \texttt{4Delta}. If \texttt{S} contains another left vertex and \texttt{J=max\{j:a\_j\ in\ S\}}, then its \texttt{a\_J} edges form a matching of size \texttt{4Delta\ J} in the right coordinates. Thus any hitting set has at least that size. Taking all original resources \texttt{x\_1,...,x\_(4Delta\ J)} hits every edge for \texttt{S}, proving

\[
\tau(E_S)=4\Delta J
                         \ge4\Delta|S|.
\]

Hence (2) holds with \texttt{r=2}. Also \texttt{Delta} is a fixed polynomial of \texttt{p}; for example \texttt{Delta\textgreater{}=p\^{}(1/3)} for all sufficiently large \texttt{q}.

\hypertarget{af-6a8212874a-a-persistent-expensive-unmatched-left-vertex}{%
\subsubsection{A persistent expensive unmatched left vertex}\label{af-6a8212874a-a-persistent-expensive-unmatched-left-vertex}}

Use the legitimate deterministic greedy order \texttt{a\_q,a\_(q-1),...,a\_1}. Within \texttt{a\_j} for \texttt{j\textgreater{}=2}, scan resources by increasing \texttt{t}. Until \texttt{a\_1} is matched, every matching edge consists of one original resource and one private resource owned by a different, already matched left vertex.

During a root-layer build, the currently forbidden original resources therefore form an initial segment \texttt{x\_1,...,x\_c}. A candidate for an unmatched \texttt{a\_j}, \texttt{j\textgreater{}=2}, is available exactly when its original resource is unmarked. Its private resource belongs to that unmatched left vertex and cannot occur in the current matching or in another left vertex's selected edge. Deleting the full matching blocker of an accepted edge adds no additional original resource, only an irrelevant private resource of a matched vertex. Thus the initial-segment description persists throughout the scan.

Suppose the unmatched set contains the entire prefix \texttt{a\_1,...,a\_L}, with \texttt{L\textgreater{}=6}, and put \texttt{j\_0=floor(L/6)}. After processing the vertices \texttt{a\_L,...,a\_(j\_0+1)}, one must have \texttt{c\textgreater{}=4Delta\ j\_0}. Otherwise, each of those \texttt{L-j\_0} vertices had more than \texttt{Delta} available original resources and accepted its full quota. This would give

\[
c\ge(L-j_0)\Delta\ge5L\Delta/6
        >4j_0\Delta,
\]

a contradiction. Any previously processed higher-index active vertices only increase \texttt{c} and cannot invalidate this argument.

Consequently every \texttt{a\_j} with \texttt{2\textless{}=j\textless{}=j\_0} has all its candidate original resources marked and accepts no edge. For \texttt{a\_1}, at most one of its \texttt{4Delta+1} clique resources remains unmarked, so it too accepts no edge. The entire prefix \texttt{a\_1,...,a\_(floor(L/6))} remains unmatched after this root layer collapses.

There is no hidden deeper operation that matches these vertices. The root layer contains at least \texttt{Delta} selected edges, while at most \texttt{q} of them have a matching blocker. Since \texttt{Delta=q\^{}3}, more than one tenth are free and the root layer collapses immediately. Lower-level swaps are not needed. Only left vertices represented among its selected edges can become newly matched in that collapse.

Starting from \texttt{L=q}, at least \texttt{Omega(log\ q)=Omega(log\ p)} root builds occur before \texttt{a\_1} can be matched. On every such build, an ordinary edge-list scan visits all \texttt{Theta(Delta\^{}2)=Theta(p)} candidate edges of \texttt{a\_1}, because none is accepted and its quota never fills. The resulting work is \texttt{Omega(p\ log\ p)}, even when edges of all inactive left vertices are skipped.

\hypertarget{af-6a8212874a-what-would-be-required-to-improve-this-cost}{%
\subsubsection{What would be required to improve this cost}\label{af-6a8212874a-what-would-be-required-to-improve-this-cost}}

The example rules out the claim that the sum of active-edge incidences scanned by these builds is automatically \texttt{O(p)}. It does not rule out a different implementation. For example, a data structure might reuse a previously discovered hitting certificate whenever the current forbidden resources still contain it. Such reuse would need a separate correctness and update-cost argument as the matching and forbidden sets change. No such general amortized bound has been established here. Accordingly the root factor in (3)--(6) is retained.
